% Pursuing Stacks (cotes 134-2 to 134-8): archive cote 134-3, site page 66 (PDF page 67), typescript page 66.
% typescript. Draft compared: ps2.tex:3410-3483.
% AI-transcribed by Claude Opus 5.5 (claude-opus-5-5, adaptive thinking, effort high), Message Batches API, batch ps1 (msgbatch_01Rv5hCaNTtUt4JHx1yqoHPL), 2026-11-12; not yet checked by eye.
\scanpage{134-3/66}
\mspage{\hand{66}}

To be completely happy, we still need a down-to-earth \xout{xrixe} sufficient
criterion to insure that $i^{*}(\hand{\Delta}_1)$ is aspheric over $e_{\unsure{A\hat{}}}$ , in the spirit \over{of}{of}
the somewhat awkward \xout{xxmdx}\inserted{[deleted]} corollary of \deleted{\unsure{p. 38.}}\inserted{end of section 36.} The following seems quite
adequate for all cases I have in mind at the present moment:

\uline{Corollary} 3. Under the general conditions of the theorem on A and \xout{xxxAx(Ca}
$i\colon A\longrightarrow$ (Cat), assume moreover we got a homotopy interval $(I,\hand{\delta}_0,\hand{\delta}_1)$ in
$A\hat{}$ \hand{*}, that A has a final object e (which is therefore also a final object
of $A\hat{}$ so we may view $\hand{\delta}_0,\hand{\delta}_1$ as maps $e\rightrightarrows I$ \deleted{\unsure{;}} i.e. elements in I(e)), and
let $i_!\colon A\hat{}\longrightarrow$ (Cat) be the canonical extension of i to $A\hat{}$(commuting with
direct limits). Assume moreover $i(e)=\hand{\Delta}_0$ (the final \xout{category}\inserted[typed]{object} in (Cat)),
and that we can find a map in (Cat)
\margin{cf def. p. 59}
\[ i_!(I)\to\hand{\Delta}_1 \]
compatible with $\hand{\delta}_0,\hand{\delta}_1$ , i.e. whose compositae with $i_!(\hand{\delta}_n)\colon\hand{\Delta}_0\to\hand{\Delta}_1$
for $n\hand{\in\{}0,1\hand{\}}$ are the two standar\over{t}{d} maps $\underline{\hand{\delta}}_0,\underline{\hand{\delta}}_1$ from $\hand{\Delta}_0$ to $\hand{\Delta}_1$ . Then the
condition b (iii) of the theorem holds, namely $i^{*}(\hand{\Delta}_1)$ is aspheric over e

Indeed, to give a map $i_!(I)\to\hand{\Delta}_1$ \inserted{in (Cat)} amounts to the same as giving
a map $I\to i^{*}(\hand{\Delta}_1)$ in $A\hat{}$ (namely $i_!$ and $i^{*}$ are adjoint), moreover the extra
condition involving $\hand{\delta}_0,\hand{\delta}_1$ just means that this map respects endpoints.
Using the composition law of intersection on $i^{*}(\hand{\Delta}_1)=(a\mapsto\mathrm{Crib}(i(a)))$,
the comparison \xout{thxx} lemma for homotopy intervals \deleted{(\unsure{p. 60}}\inserted{(section 37)}) \inserted{,} implies that
$i^{*}(\hand{\Delta}_1)$ is aspheric over e , qed.

In the cases I have in mind, I is even an element of A hence
$i_!(I)=i(I)$, moreover $i(I)=\hand{\Delta}_1$ and the map $i_!(I)\to\hand{\Delta}_1$ above is the
\uline{identity}\over{.}{!} The choice of the functor i is in every case ``the most natural
one'' (discarding however the clumsy $i_0=i_A$ , and trying to get away
with categories i(a) which \deleted{are finite} give the simplest imaginable descrip-%
tion of objects and arrows of the would-be test category A), and the choi-%
ce of I itself is still more evident - it is \uline{the} object of A (or \uline{one} among
the objects, in c ases such as $A=\hand{\Delta}^n$ , giving rise to simplicial \uline{multi}com
plexes ...) which suggests most strongly the picture of an ``interval''. Thus
\margin{\unsure{$\updownarrow$}}the one key verification we are left with (all the rest being ``formal'' in
terms of what precedes) is the asphericity of I over e, i.e. that all
the products Ixa are aspher\deleted{\unsure{r}}ic.

\margin{\unsure{$\updownarrow$}}39. Maybe it is time now to come back to the property of total asphericity
of $A\hat{}$, expressed by the condition T 2) on A , namely that the product in
$A\hat{}$of any two elements in A is aspheric. As we saw, \inserted{[if $A\neq\emptyset$],} this implies already
asphericity of A i.e. of \xout{thx} the \unsure{topos} $A\hat{}$ . In our present setting, total
asphericity is of special interest only when coupled with the property
T 3 , which amounts to saying that the Lawvere element $L_A$ in $A\hat{}$ is
